% Options for packages loaded elsewhere
% Options for packages loaded elsewhere
\PassOptionsToPackage{unicode}{hyperref}
\PassOptionsToPackage{hyphens}{url}
%
\documentclass[
  11pt,
  ignorenonframetext,
  aspectratio=169,
]{beamer}
\newif\ifbibliography
\usepackage{pgfpages}
\setbeamertemplate{caption}[numbered]
\setbeamertemplate{caption label separator}{: }
\setbeamercolor{caption name}{fg=normal text.fg}
\beamertemplatenavigationsymbolsempty
% remove section numbering
\setbeamertemplate{part page}{
  \centering
  \begin{beamercolorbox}[sep=16pt,center]{part title}
    \usebeamerfont{part title}\insertpart\par
  \end{beamercolorbox}
}
\setbeamertemplate{section page}{
  \centering
  \begin{beamercolorbox}[sep=12pt,center]{section title}
    \usebeamerfont{section title}\insertsection\par
  \end{beamercolorbox}
}
\setbeamertemplate{subsection page}{
  \centering
  \begin{beamercolorbox}[sep=8pt,center]{subsection title}
    \usebeamerfont{subsection title}\insertsubsection\par
  \end{beamercolorbox}
}
% Prevent slide breaks in the middle of a paragraph
\widowpenalties 1 10000
\raggedbottom
\AtBeginPart{
  \frame{\partpage}
}
\AtBeginSection{
  \ifbibliography
  \else
    \frame{\sectionpage}
  \fi
}
\AtBeginSubsection{
  \frame{\subsectionpage}
}
\usepackage{iftex}
\ifPDFTeX
  \usepackage[T1]{fontenc}
  \usepackage[utf8]{inputenc}
  \usepackage{textcomp} % provide euro and other symbols
\else % if luatex or xetex
  \usepackage{unicode-math} % this also loads fontspec
  \defaultfontfeatures{Scale=MatchLowercase}
  \defaultfontfeatures[\rmfamily]{Ligatures=TeX,Scale=1}
\fi
\usepackage{lmodern}

\ifPDFTeX\else
  % xetex/luatex font selection
\fi
% Use upquote if available, for straight quotes in verbatim environments
\IfFileExists{upquote.sty}{\usepackage{upquote}}{}
\IfFileExists{microtype.sty}{% use microtype if available
  \usepackage[]{microtype}
  \UseMicrotypeSet[protrusion]{basicmath} % disable protrusion for tt fonts
}{}
\makeatletter
\@ifundefined{KOMAClassName}{% if non-KOMA class
  \IfFileExists{parskip.sty}{%
    \usepackage{parskip}
  }{% else
    \setlength{\parindent}{0pt}
    \setlength{\parskip}{6pt plus 2pt minus 1pt}}
}{% if KOMA class
  \KOMAoptions{parskip=half}}
\makeatother

\usepackage{color}
\usepackage{fancyvrb}
\newcommand{\VerbBar}{|}
\newcommand{\VERB}{\Verb[commandchars=\\\{\}]}
\DefineVerbatimEnvironment{Highlighting}{Verbatim}{commandchars=\\\{\}}
% Add ',fontsize=\small' for more characters per line
\usepackage{framed}
\definecolor{shadecolor}{RGB}{241,243,245}
\newenvironment{Shaded}{\begin{snugshade}}{\end{snugshade}}
\newcommand{\AlertTok}[1]{\textcolor[rgb]{0.68,0.00,0.00}{#1}}
\newcommand{\AnnotationTok}[1]{\textcolor[rgb]{0.37,0.37,0.37}{#1}}
\newcommand{\AttributeTok}[1]{\textcolor[rgb]{0.40,0.45,0.13}{#1}}
\newcommand{\BaseNTok}[1]{\textcolor[rgb]{0.68,0.00,0.00}{#1}}
\newcommand{\BuiltInTok}[1]{\textcolor[rgb]{0.00,0.23,0.31}{#1}}
\newcommand{\CharTok}[1]{\textcolor[rgb]{0.13,0.47,0.30}{#1}}
\newcommand{\CommentTok}[1]{\textcolor[rgb]{0.37,0.37,0.37}{#1}}
\newcommand{\CommentVarTok}[1]{\textcolor[rgb]{0.37,0.37,0.37}{\textit{#1}}}
\newcommand{\ConstantTok}[1]{\textcolor[rgb]{0.56,0.35,0.01}{#1}}
\newcommand{\ControlFlowTok}[1]{\textcolor[rgb]{0.00,0.23,0.31}{\textbf{#1}}}
\newcommand{\DataTypeTok}[1]{\textcolor[rgb]{0.68,0.00,0.00}{#1}}
\newcommand{\DecValTok}[1]{\textcolor[rgb]{0.68,0.00,0.00}{#1}}
\newcommand{\DocumentationTok}[1]{\textcolor[rgb]{0.37,0.37,0.37}{\textit{#1}}}
\newcommand{\ErrorTok}[1]{\textcolor[rgb]{0.68,0.00,0.00}{#1}}
\newcommand{\ExtensionTok}[1]{\textcolor[rgb]{0.00,0.23,0.31}{#1}}
\newcommand{\FloatTok}[1]{\textcolor[rgb]{0.68,0.00,0.00}{#1}}
\newcommand{\FunctionTok}[1]{\textcolor[rgb]{0.28,0.35,0.67}{#1}}
\newcommand{\ImportTok}[1]{\textcolor[rgb]{0.00,0.46,0.62}{#1}}
\newcommand{\InformationTok}[1]{\textcolor[rgb]{0.37,0.37,0.37}{#1}}
\newcommand{\KeywordTok}[1]{\textcolor[rgb]{0.00,0.23,0.31}{\textbf{#1}}}
\newcommand{\NormalTok}[1]{\textcolor[rgb]{0.00,0.23,0.31}{#1}}
\newcommand{\OperatorTok}[1]{\textcolor[rgb]{0.37,0.37,0.37}{#1}}
\newcommand{\OtherTok}[1]{\textcolor[rgb]{0.00,0.23,0.31}{#1}}
\newcommand{\PreprocessorTok}[1]{\textcolor[rgb]{0.68,0.00,0.00}{#1}}
\newcommand{\RegionMarkerTok}[1]{\textcolor[rgb]{0.00,0.23,0.31}{#1}}
\newcommand{\SpecialCharTok}[1]{\textcolor[rgb]{0.37,0.37,0.37}{#1}}
\newcommand{\SpecialStringTok}[1]{\textcolor[rgb]{0.13,0.47,0.30}{#1}}
\newcommand{\StringTok}[1]{\textcolor[rgb]{0.13,0.47,0.30}{#1}}
\newcommand{\VariableTok}[1]{\textcolor[rgb]{0.07,0.07,0.07}{#1}}
\newcommand{\VerbatimStringTok}[1]{\textcolor[rgb]{0.13,0.47,0.30}{#1}}
\newcommand{\WarningTok}[1]{\textcolor[rgb]{0.37,0.37,0.37}{\textit{#1}}}

\usepackage{longtable,booktabs,array}
\newcounter{none} % for unnumbered tables
\usepackage{calc} % for calculating minipage widths
\usepackage{caption}
% Make caption package work with longtable
\makeatletter
\def\fnum@table{\tablename~\thetable}
\makeatother
\usepackage{graphicx}
\makeatletter
\newsavebox\pandoc@box
\newcommand*\pandocbounded[1]{% scales image to fit in text height/width
  \sbox\pandoc@box{#1}%
  \Gscale@div\@tempa{\textheight}{\dimexpr\ht\pandoc@box+\dp\pandoc@box\relax}%
  \Gscale@div\@tempb{\linewidth}{\wd\pandoc@box}%
  \ifdim\@tempb\p@<\@tempa\p@\let\@tempa\@tempb\fi% select the smaller of both
  \ifdim\@tempa\p@<\p@\scalebox{\@tempa}{\usebox\pandoc@box}%
  \else\usebox{\pandoc@box}%
  \fi%
}
% Set default figure placement to htbp
\def\fps@figure{htbp}
\makeatother





\setlength{\emergencystretch}{3em} % prevent overfull lines

\providecommand{\tightlist}{%
  \setlength{\itemsep}{0pt}\setlength{\parskip}{0pt}}



 


\setbeamertemplate{navigation symbols}{}
\setbeamertemplate{footline}[frame number]
\setbeamertemplate{itemize items}[circle]
\setbeamertemplate{itemize subitem}[triangle]
\setbeamercolor{frametitle}{fg=black}
\setbeamerfont{frametitle}{series=\bfseries}
\usepackage{amsmath}
\usepackage{booktabs}
\usepackage{array}
\newcolumntype{R}[1]{>{\raggedright\arraybackslash}p{#1}}
\usepackage{fancyvrb}
\fvset{fontsize=\scriptsize}
\makeatletter
\def\verbatim@font{\scriptsize\ttfamily}
\makeatother
\makeatletter
\@ifpackageloaded{caption}{}{\usepackage{caption}}
\AtBeginDocument{%
\ifdefined\contentsname
  \renewcommand*\contentsname{Table of contents}
\else
  \newcommand\contentsname{Table of contents}
\fi
\ifdefined\listfigurename
  \renewcommand*\listfigurename{List of Figures}
\else
  \newcommand\listfigurename{List of Figures}
\fi
\ifdefined\listtablename
  \renewcommand*\listtablename{List of Tables}
\else
  \newcommand\listtablename{List of Tables}
\fi
\ifdefined\figurename
  \renewcommand*\figurename{Figure}
\else
  \newcommand\figurename{Figure}
\fi
\ifdefined\tablename
  \renewcommand*\tablename{Table}
\else
  \newcommand\tablename{Table}
\fi
}
\@ifpackageloaded{float}{}{\usepackage{float}}
\floatstyle{ruled}
\@ifundefined{c@chapter}{\newfloat{codelisting}{h}{lop}}{\newfloat{codelisting}{h}{lop}[chapter]}
\floatname{codelisting}{Listing}
\newcommand*\listoflistings{\listof{codelisting}{List of Listings}}
\makeatother
\makeatletter
\makeatother
\makeatletter
\@ifpackageloaded{caption}{}{\usepackage{caption}}
\@ifpackageloaded{subcaption}{}{\usepackage{subcaption}}
\makeatother

\usepackage{bookmark}
\IfFileExists{xurl.sty}{\usepackage{xurl}}{} % add URL line breaks if available
\urlstyle{same}
\hypersetup{
  pdftitle={Dynamic Panels},
  pdfauthor={Professor Ji-Woong Chung},
  hidelinks,
  pdfcreator={LaTeX via pandoc}}


\title{\texorpdfstring{Dynamic Panels}{Dynamic Panels}}
\subtitle{\texorpdfstring{BUSS975: Causal Inference in Financial
Research}{BUSS975: Causal Inference in Financial Research}}
\author{Professor Ji-Woong Chung}
\date{}
\institute{Korea University}

\begin{document}
\frame{\titlepage}


\begin{frame}{Where we left off}
\protect\phantomsection\label{where-we-left-off}
\begin{itemize}
\tightlist
\item
  The panel lecture rested on one claim: anything about a firm that
  \textbf{does not change} cannot explain why its outcome changed.
  Demean, and the firm effect is gone.
\item
  The price was \textbf{strict exogeneity}: today's shock must not feed
  into tomorrow's regressors.
\end{itemize}

\vspace{0.4em}

\begin{itemize}
\tightlist
\item
  Finance is full of outcomes that depend on \textbf{their own past}:
  leverage, payout, cash, investment.
\item
  Put yesterday's \(Y\) on the right-hand side and strict exogeneity is
  not merely doubtful. \textbf{It is false by construction.}
\item
  This lecture: what breaks, how badly, how to fix it --- and when the
  fix does not work.
\end{itemize}
\end{frame}

\begin{frame}{The plan}
\protect\phantomsection\label{the-plan}
\begin{enumerate}
\tightlist
\item
  \textbf{Why the past belongs on the right-hand side} --- and one rule
  that settles every calculation.
\item
  \textbf{Two estimators, two wrong answers.} They err in opposite
  directions, which turns out to be useful.
\item
  \textbf{The fix: instruments from the past.} Anderson--Hsiao, then
  Arellano--Bond.
\item
  \textbf{Why the fix can fail} --- and system GMM, which fails
  differently.
\item
  \textbf{In practice:} a case study from capital structure, the code,
  and a checklist.
\end{enumerate}

\vspace{0.4em}

\begin{itemize}
\tightlist
\item
  Two running questions. \emph{Which comparison is this estimator
  making?} And, new this week: \emph{which way is it wrong?}
\end{itemize}
\end{frame}

\section{1. Why the Past Belongs on the Right-Hand
Side}\label{why-the-past-belongs-on-the-right-hand-side}

\begin{frame}{Most things in finance drift}
\protect\phantomsection\label{most-things-in-finance-drift}
They do not jump to a target. They close part of the gap each year.

\begin{center}\small
\begin{tabular}{ll}
\toprule
Leverage & firms close only part of the gap to target each year \\
Payout & managers smooth dividends, and cut only as a last resort \\
Cash & drawn down quickly, rebuilt slowly \\
\bottomrule
\end{tabular}
\end{center}

\[
Y_{it} \;=\; \rho\,Y_{i,t-1} + \delta D_{it} + u_i + \varepsilon_{it}, \qquad |\rho| < 1
\]

\begin{itemize}
\tightlist
\item
  \(\rho\) \textbf{is the economics.} \(1-\rho\) is the share of the gap
  closed per year; a shock has half-life \(\ln(0.5)/\ln\rho\). At
  \(\rho = 0.6\) that is \(1.4\) years; at \(\rho = 0.9\), \(6.6\)
  years. \textbf{Different theories of the firm.}
\item
  \(\delta\) is now the \textbf{impact} effect. The long-run effect is
  \(\delta/(1-\rho)\) --- get \(\rho\) wrong and you get \emph{that}
  wrong too.
\end{itemize}
\end{frame}

\begin{frame}{What a firm effect does to a series}
\protect\phantomsection\label{what-a-firm-effect-does-to-a-series}
\begin{center}
\includegraphics[width=0.78\linewidth,height=\textheight,keepaspectratio]{06b_dynamic_panel_files/figure-beamer/two-firms-1.pdf}
\end{center}

Two firms, same \(\rho = 0.6\), same shocks in distribution ---
different \(u_i\). Back-substitute the model into itself: \[
Y_{it} \;=\; \frac{u_i}{1-\rho} \;+\; \delta\sum_{j\ge0}\rho^{\,j}D_{i,t-j} \;+\; \sum_{j\ge0}\rho^{\,j}\varepsilon_{i,t-j}
\]

\begin{itemize}
\tightlist
\item
  The firm effect shifts the \textbf{level} --- not the volatility, not
  the persistence. Each firm wiggles around its own long-run level
  \(u_i/(1-\rho)\).
\item
  \textbf{Remember this picture.} An estimator that cannot tell
  \emph{high because it was high last year} from \emph{high because this
  firm is simply high} will get \(\rho\) wrong.
\end{itemize}
\end{frame}

\begin{frame}{One rule settles every calculation}
\protect\phantomsection\label{one-rule-settles-every-calculation}
\(Y_{is}\) is assembled from \(u_i\) and from shocks \textbf{dated \(s\)
and earlier} --- nothing later. If \(\varepsilon_{it}\) is serially
uncorrelated and unrelated to \(u_i\): \[
\boxed{\;\mathbb{E}[Y_{is}\,\varepsilon_{it}] = 0 \;\;\text{whenever}\;\; t > s,
\qquad
\mathbb{E}[Y_{is}\,\varepsilon_{is}] = \sigma^2_\varepsilon\;}
\]

\begin{itemize}
\tightlist
\item
  In words: \textbf{a later shock has nothing to do with an earlier
  \(Y\)}; the same-date shock enters \(Y\) with coefficient one.
\item
  Every derivation in this lecture is a \textbf{date check}: write out a
  product, compare the two dates, and it is either zero or
  \(\sigma^2_\varepsilon\).
\item
  One clause does quiet work: \emph{serially uncorrelated}. It turns out
  to be the identifying assumption of the whole enterprise --- and it is
  testable.
\end{itemize}
\end{frame}

\begin{frame}{The problem you cannot assume away}
\protect\phantomsection\label{the-problem-you-cannot-assume-away}
Lag the model one period and the difficulty is on the page: \[
Y_{i,t-1} \;=\; \rho\,Y_{i,t-2} + \delta D_{i,t-1} + \underbrace{u_i}_{\text{\textbf{also in today's error}}} + \;\varepsilon_{i,t-1}
\]

\begin{itemize}
\tightlist
\item
  Every endogenous regressor so far was endogenous because of something
  we \textbf{failed to measure}. This one is endogenous \textbf{by
  construction}.
\item
  \(Y_{i,t-1}\) contains \(u_i\), and so does the error: \textbf{pooled
  OLS is biased.}
\item
  \(\varepsilon_{it}\) moves \(Y_{it}\), which \emph{is} next year's
  regressor: strict exogeneity is false, so \textbf{fixed effects is
  biased too.}
\item
  The useful part: \textbf{they fail in opposite directions.}
\end{itemize}
\end{frame}

\section{2. Two Estimators, Two Wrong
Answers}\label{two-estimators-two-wrong-answers}

\begin{frame}{Why OLS reads heterogeneity as persistence}
\protect\phantomsection\label{why-ols-reads-heterogeneity-as-persistence}
Strip the model to the bone. Set \(\rho = 0\) --- \textbf{no persistence
whatsoever}, only permanent differences between firms: \[
Y_{it} \;=\; u_i + \varepsilon_{it}
\]

\begin{itemize}
\tightlist
\item
  A high-\(u_i\) firm is high this year \textbf{and} was high last year
  --- not because the past caused the present, but because the firm is
  simply like that.
\item
  Pooled OLS sees \(Y\) following \(Y\) and calls it persistence: \[
  \operatorname{plim} \hat\rho_{\text{OLS}} \;=\; \frac{\operatorname{Cov}(Y_{it}, Y_{i,t-1})}{\operatorname{Var}(Y_{i,t-1})} \;=\; \frac{\sigma^2_u}{\sigma^2_u + \sigma^2_\varepsilon} \;>\; 0 \;=\; \rho
  \]
\item
  With \(\sigma_u = \sigma_\varepsilon = 1\) that is \(0.5\). Simulated
  (\(T=5\), \(100{,}000\) firms):
  \(\hat\rho_{\text{OLS}} = \mathbf{0.501}\) when the truth is
  \textbf{zero}.
\item
  \textbf{More years do not help.} More years of the same firm being
  itself is more of the same evidence.
\end{itemize}
\end{frame}

\begin{frame}{Why demeaning manufactures mean reversion}
\protect\phantomsection\label{why-demeaning-manufactures-mean-reversion}
\textbf{The story.} A lucky year raises \(Y_{i,t-1}\) \textbf{and} the
firm's average. Measured against that inflated average, the firm's other
years look low --- so the firm \emph{seems} to revert even when nothing
reverts.

\vspace{0.3em}

In symbols: the firm mean contains the errors, so after demeaning \[
\underbrace{Y_{i,t-1} - \bar Y_{i,-1}}_{\text{regressor}}
\quad\text{is correlated with}\quad
\underbrace{\varepsilon_{it} - \bar\varepsilon_{i}}_{\text{error}}:
\qquad
\operatorname{Cov}\big(Y_{i,t-1},\, -\bar\varepsilon_i\big) \;<\; 0, \;\; \text{of order } \sigma^2_\varepsilon/T
\]

\begin{itemize}
\tightlist
\item
  So \(\hat\rho_{\text{FE}}\) is biased \textbf{down}: \textbf{Nickell
  (1981) bias.}
\item
  The same \(\rho = 0\) experiment with fixed effects:
  \(\hat\rho_{\text{FE}} = \mathbf{-0.251}\), almost exactly
  \(-1/(T-1)\). OLS invented persistence; demeaning invents \textbf{mean
  reversion}.
\item
  It happens \textbf{even if \(u_i\) does not exist} --- demeaning alone
  does it --- and it fades only at rate \(1/T\).
\end{itemize}
\end{frame}

\begin{frame}{A derivation you can do at the board: \(T = 3\)}
\protect\phantomsection\label{a-derivation-you-can-do-at-the-board-t-3}
With \(T = 3\) there is \textbf{one differenced observation per firm},
and \(\hat\rho_{\text{FE}}\) is exactly least squares on
\(\Delta Y_{i3} = \rho\,\Delta Y_{i2} + \Delta\varepsilon_{i3}\): \[
\operatorname{plim}(\hat\rho_{\text{FE}} - \rho) \;=\; \frac{\mathbb{E}[\Delta Y_{i2}\,\Delta\varepsilon_{i3}]}{\mathbb{E}[(\Delta Y_{i2})^2]}
\]

\begin{itemize}
\tightlist
\item
  \textbf{Numerator, by the rule.} Expand
  \((Y_{i2} - Y_{i1})(\varepsilon_{i3} - \varepsilon_{i2})\) into four
  products. Three pair a \(Y\) with a \emph{later} shock: zero. One
  pairs \(Y_{i2}\) with its \emph{own} shock, with a minus sign:
  \(-\sigma^2_\varepsilon\).
\item
  \textbf{Denominator, from the picture.} \(u_i\) cancels in a
  difference; what is left has variance
  \(2\sigma^2_\varepsilon/(1+\rho)\). \[
  \operatorname{plim}(\hat\rho_{\text{FE}} - \rho) \;=\; \frac{-\sigma^2_\varepsilon}{2\sigma^2_\varepsilon/(1+\rho)} \;=\; -\,\frac{1+\rho}{2}
  \]
\item
  \textbf{The plim of \(\hat\rho_{\text{FE}}\) is negative for every
  \(\rho<1\)}: leverage that is genuinely \(90\%\) persistent is
  reported as mean-reverting. (Hansen 2022, §17.37.)
\item
  Keep that \(-\sigma^2_\varepsilon\) in mind. It comes back in Section
  3.
\end{itemize}
\end{frame}

\begin{frame}{Watching both biases at once}
\protect\phantomsection\label{watching-both-biases-at-once}
Simulate \(Y_{it} = \rho Y_{i,t-1} + u_i + \varepsilon_{it}\) with
\(\rho = 0.6\) and \(4{,}000\) firms, varying the number of years \(T\).

\begin{center}\small
\begin{tabular}{lcccc}
\toprule
Years per firm & $T=3$ & $T=5$ & $T=10$ & $T=30$ \\
\midrule
Pooled OLS \quad(biased \textbf{up}) & $0.919$ & $0.918$ & $0.919$ & $0.918$ \\
Fixed effects \quad(biased \textbf{down}) & $-0.189$ & $0.147$ & $0.396$ & $0.543$ \\
\quad \emph{Nickell (1981) plim, for comparison} & \itshape -0.200 & \itshape 0.153 & \itshape 0.401 & \itshape 0.542 \\
\midrule
Truth & $0.600$ & $0.600$ & $0.600$ & $0.600$ \\
\bottomrule
\end{tabular}
\end{center}

\begin{itemize}
\tightlist
\item
  \textbf{The OLS row never moves.} Heterogeneity is not resolved by
  more years.
\item
  \textbf{At \(T = 3\) fixed effects is negative} for a process with
  persistence \(+0.6\) --- as \(-(1+\rho)/2\) predicted.
\item
  The FE row tracks Nickell's plim throughout: this is bias, not noise.
\end{itemize}
\end{frame}

\begin{frame}{Does it fade with \(T\)? Slowly.}
\protect\phantomsection\label{does-it-fade-with-t-slowly.}
\begin{center}
\includegraphics[width=0.85\linewidth,height=\textheight,keepaspectratio]{06b_dynamic_panel_files/figure-beamer/nickell-plot-1.pdf}
\end{center}

\begin{itemize}
\tightlist
\item
  The bias is \(O(1/T)\): doubling the panel halves it. Finance panels
  are typically \(10\)--\(20\) years, often unbalanced.
\item
  At \(T = 30\) the bias is still \(-0.058\) at \(\rho = 0.6\) and
  \(-0.080\) at \(\rho = 0.9\). \textbf{The more persistent the series,
  the slower the bias leaves.}
\end{itemize}
\end{frame}

\begin{frame}{Two wrongs that bound the right answer}
\protect\phantomsection\label{two-wrongs-that-bound-the-right-answer}
Neither estimator is consistent. But you know the \textbf{sign of each
bias}, so together they fence the truth in: \[
\hat\rho_{\text{FE}} \;<\; \rho \;<\; \hat\rho_{\text{OLS}}
\]

\begin{itemize}
\tightlist
\item
  In the table at \(T=5\): \(0.15 < 0.60 < 0.92\). Wide, and
  \textbf{honest}. A single number from either estimator would not be.
\item
  \textbf{The standard first move} (Bond 2002): run OLS and fixed
  effects \emph{before} anything clever, and keep both numbers. They are
  the goalposts.
\item
  \textbf{Any estimator you trust must land between them.} One that
  lands outside tells you something is wrong before you interpret a
  single coefficient.
\item
  A bracket is a statement about \textbf{bias}, not a confidence
  interval.
\end{itemize}

\vspace{0.2em}
\footnotesize

(A second bracket, for \(\delta\) when you are unsure whether the
lagged-\(Y\) or the fixed-effects model is right, is in the backup.)
\normalsize
\end{frame}

\section{3. The Fix: Instruments from the
Past}\label{the-fix-instruments-from-the-past}

\begin{frame}{The fix, in one idea: Anderson--Hsiao}
\protect\phantomsection\label{the-fix-in-one-idea-andersonhsiao}
Difference the model to kill \(u_i\): \[
\Delta Y_{it} \;=\; \rho\,\Delta Y_{i,t-1} + \delta\,\Delta D_{it} + \Delta\varepsilon_{it}
\]

\begin{itemize}
\tightlist
\item
  But \(\Delta Y_{i,t-1}\) contains \(\varepsilon_{i,t-1}\), and so does
  \(\Delta\varepsilon_{it}\): the same \(-\sigma^2_\varepsilon\) as the
  \(T=3\) derivation. \textbf{Differencing is not the fix; it is where
  the fixed-effects bias comes from.}
\item
  \textbf{So go back one more lag.} \(Y_{i,t-2}\) is \textbf{far enough
  back to be clean}, yet \textbf{close enough to still predict}
  \(\Delta Y_{i,t-1}\).
\item
  That is \textbf{Anderson--Hsiao (1982)}: IV on the differenced
  equation, instrumenting \(\Delta Y_{i,t-1}\) with \(Y_{i,t-2}\). It
  needs \(T \ge 3\).
\item
  The lagged outcome caused the problem, and a \emph{longer} lag of the
  same variable solves it --- no new data, only the history you already
  have.
\end{itemize}
\end{frame}

\begin{frame}{One date apart: a broken regressor and a valid instrument}
\protect\phantomsection\label{one-date-apart-a-broken-regressor-and-a-valid-instrument}
Check each piece against the rule, with
\(\Delta\varepsilon_{it} = \varepsilon_{it} - \varepsilon_{i,t-1}\):

\begin{center}\small
\begin{tabular}{lccc}
\toprule
 & $\times\,\varepsilon_{it}$ & $\times\,\varepsilon_{i,t-1}$ & $\mathbb{E}[\,\cdot\,\Delta\varepsilon_{it}]$ \\
\midrule
$Y_{i,t-1}$ & $0$ (later shock) & $\boldsymbol{\sigma^2_\varepsilon}$ \textbf{(same date)} & $-\sigma^2_\varepsilon$ \\
$Y_{i,t-2}$ & $0$ (later shock) & $0$ (later shock) & $0$ \\
\midrule
$\Delta Y_{i,t-1} = Y_{i,t-1} - Y_{i,t-2}$ & & & $\boldsymbol{-\sigma^2_\varepsilon}$ \\
\bottomrule
\end{tabular}
\end{center}

\begin{itemize}
\tightlist
\item
  \textbf{Validity.} The regressor inherits the one same-date term; the
  instrument, one year older, has none. \textbf{The difference between a
  broken regressor and a valid instrument is one date.}
\item
  \textbf{Relevance.} \(u_i\) cancels inside \(\Delta Y_{i,t-1}\),
  leaving
  \(\operatorname{Cov}(Y_{i,t-2},\Delta Y_{i,t-1}) = -\sigma^2_\varepsilon/(1+\rho) \neq 0\).
  Nonzero --- but Section 4 shows it can still be weak.
\item
  \textbf{The zero row rests on the quiet clause.} If \(\varepsilon\) is
  serially correlated (fit AR(1) when the truth is AR(2)),
  \(\varepsilon_{i,t-2}\) inside \(Y_{i,t-2}\) predicts
  \(\varepsilon_{i,t-1}\) and the instrument is invalid. \textbf{Lag
  length is the identifying assumption} --- tested two slides on.
\end{itemize}
\end{frame}

\begin{frame}{Arellano--Bond: use every lag you have}
\protect\phantomsection\label{arellanobond-use-every-lag-you-have}
Anderson--Hsiao uses one moment condition out of many. At date \(t\),
\textbf{every} \(Y\) dated \(t-2\) or earlier passes the date check
(Holtz-Eakin, Newey and Rosen 1988; Arellano and Bond 1991):

\begin{center}\small
\begin{tabular}{lll}
\toprule
Differenced equation at & valid instruments for $\Delta Y_{i,t-1}$ & how many \\
\midrule
$t = 3$ & $Y_{i1}$ & 1 \\
$t = 4$ & $Y_{i1},\, Y_{i2}$ & 2 \\
$t = 5$ & $Y_{i1},\, Y_{i2},\, Y_{i3}$ & 3 \\
$\;\;\vdots$ & & \\
$t = T$ & $Y_{i1},\, \dots,\, Y_{i,T-2}$ & $T-2$ \\
\bottomrule
\end{tabular}
\end{center}

\begin{itemize}
\tightlist
\item
  A \textbf{block-diagonal} instrument matrix with \((T-1)(T-2)/2\)
  columns; overidentified once \(T > 3\), so the estimator is
  \textbf{GMM}.
\item
  \textbf{Gain:} efficiency --- one-instrument Anderson--Hsiao is very
  noisy. \textbf{Cost:} the count grows with \(T^2\): \(91\) instruments
  at \(T = 15\), \(406\) at \(T = 30\).
\end{itemize}
\end{frame}

\begin{frame}{Too many instruments}
\protect\phantomsection\label{too-many-instruments}
Roodman (2009): GMM with many instruments relative to \(N\) is quietly a
bad estimator.

\begin{itemize}
\tightlist
\item
  \textbf{It overfits.} With enough instruments the first stage
  reproduces \(\Delta Y_{i,t-1}\) almost exactly, contaminated part
  included, and GMM slides back toward \textbf{OLS on its own equation}.
  For difference GMM that is OLS on the differences --- biased
  \textbf{down}, below even fixed effects. (For system GMM, Section 4,
  it is pooled OLS: biased \textbf{up}.)
\item
  \textbf{It disarms the overidentification test.} Hansen's \(J\)
  becomes too weak to reject anything; a \(p\)-value of \(0.99\) is a
  warning sign, not a clean bill of health.
\item
  \textbf{Two standard remedies.} \emph{Cap the lag depth} (say
  \(Y_{i,t-2}\) to \(Y_{i,t-4}\)). \emph{Collapse}: one column per lag,
  so the count grows with \(T\), not \(T^2\).
\item
  Keep the count \textbf{well below \(N\)}, and \textbf{report it}.
\end{itemize}
\end{frame}

\begin{frame}[fragile]{Two-step weights, and the tests you must report}
\protect\phantomsection\label{two-step-weights-and-the-tests-you-must-report}
\begin{itemize}
\tightlist
\item
  \textbf{One step or two.} Two-step GMM re-weights with first-step
  residuals: efficient, but its standard errors are \textbf{badly biased
  down} in finite samples. Use the \textbf{Windmeijer (2005)} correction
  (\texttt{twostep\ vce(robust)} in Stata). If one-step robust and
  two-step Windmeijer disagree, something is weak.
\item
  \textbf{AR(1) in the differenced residuals --- expected.}
  \(\Delta\varepsilon_{it}\) and \(\Delta\varepsilon_{i,t-1}\) share
  \(\varepsilon_{i,t-1}\).
\item
  \textbf{AR(2) --- must be absent.} \(\Delta\varepsilon_{it}\) and
  \(\Delta\varepsilon_{i,t-2}\) share nothing \emph{unless
  \(\varepsilon\) is serially correlated}. This is the test of the quiet
  clause; fail it and \textbf{every instrument is invalid}.
\item
  \textbf{Hansen \(J\)}, with the instrument count beside it.
\item
  \textbf{The bracket.} Difference GMM near or below
  \(\hat\rho_{\text{FE}}\): weak or too many instruments --- they pull
  the same way. Anything above \(\hat\rho_{\text{OLS}}\): a violated
  assumption (Section 4 shows one).
\end{itemize}
\end{frame}

\section{4. Why the Fix Can Fail}\label{why-the-fix-can-fail}

\begin{frame}{Weak instruments: the first stage, in closed form}
\protect\phantomsection\label{weak-instruments-the-first-stage-in-closed-form}
Blundell and Bond (1998) wrote the first stage,
\(\Delta Y_{i,t-1} = \gamma\,Y_{i,t-2} + v_{it}\), and divided
covariance by variance: \[
\begin{aligned}
\gamma \;&=\; \frac{\operatorname{Cov}(Y_{i,t-2},\Delta Y_{i,t-1})}{\operatorname{Var}(Y_{i,t-2})}
\;=\; \frac{-\sigma^2_\varepsilon/(1+\rho)}{\sigma^2_u/(1-\rho)^2 + \sigma^2_\varepsilon/(1-\rho^2)} \\[0.3em]
&=\; (\rho - 1)\,\frac{k}{k + \sigma^2_u/\sigma^2_\varepsilon},
\qquad k = \frac{1-\rho}{1+\rho}
\end{aligned}
\]

\vspace{-0.4em}

\begin{itemize}
\tightlist
\item
  The covariance on top never vanishes; \textbf{the damage is all in the
  denominator.} The \emph{level} \(Y_{i,t-2}\) is swamped by things that
  say nothing about the \emph{change}.
\item
  \textbf{As \(\rho \to 1\)}, a random walk's level says nothing about
  its next step.
\item
  \textbf{As \(\sigma^2_u/\sigma^2_\varepsilon\) grows}, the level is
  mostly the firm's permanent height, which cancels out of the change.
  In the two-firm picture: the heights tell you nothing about the
  wiggles.
\item
  Both are the normal case in corporate finance: \textbf{persistent
  series, large firm effects.}
\end{itemize}
\end{frame}

\begin{frame}{A picture of the first stage}
\protect\phantomsection\label{a-picture-of-the-first-stage}
\begin{center}
\includegraphics[width=0.85\linewidth,height=\textheight,keepaspectratio]{06b_dynamic_panel_files/figure-beamer/gamma-plot-1.pdf}
\end{center}

\begin{itemize}
\tightlist
\item
  At \(\rho = 0.6\) with equal variances, \(\gamma = -0.08\): a usable
  first stage. At \(\rho = 0.95\), \(\gamma = -0.0013\) --- essentially
  nothing.
\item
  Leverage, payout and cash all sit in the right-hand third of this
  plot. \textbf{The instrument is weakest exactly where you most need
  it.}
\end{itemize}
\end{frame}

\begin{frame}{Watching it fail}
\protect\phantomsection\label{watching-it-fail}
Same design: \(500\) firms, \(T = 6\), \(200\) replications.
\textbf{Median} and \textbf{10th--90th percentile} across replications.

\begin{center}\small
\begin{tabular}{lcc}
\toprule
 & $\rho = 0.6$ & $\rho = 0.95$ \\
\midrule
Pooled OLS & $0.92$ \;{\footnotesize[$0.91$, $0.93$]} & $1.00$ \;{\footnotesize[$1.00$, $1.00$]} \\
Fixed effects & $0.24$ \;{\footnotesize[$0.21$, $0.26$]} & $0.47$ \;{\footnotesize[$0.44$, $0.49$]} \\
Anderson--Hsiao (one instrument) & $0.59$ \;{\footnotesize[$0.50$, $0.71$]} & $0.88$ \;{\footnotesize[$-0.55$, $2.14$]} \\
Arellano--Bond (all lags, one-step) & $0.58$ \;{\footnotesize[$0.50$, $0.67$]} & $0.60$ \;{\footnotesize[$0.16$, $0.93$]} \\
\bottomrule
\end{tabular}
\end{center}

\begin{itemize}
\tightlist
\item
  \textbf{At \(\rho = 0.6\) the fix works.} Both IV estimators sit
  inside the bracket, near the truth; Arellano--Bond is the tighter.
\item
  \textbf{At \(\rho = 0.95\) it does not.} Anderson--Hsiao is scattered
  across the number line; Arellano--Bond is pulled far down with an
  enormous spread --- the signature of weak instruments.
\end{itemize}
\end{frame}

\begin{frame}{System GMM: the idea, and its price}
\protect\phantomsection\label{system-gmm-the-idea-and-its-price}
Blundell and Bond's repair: keep the differenced equation, and
\textbf{add the levels equation back}, instrumented the other way round:
\[
Y_{it} = \rho\,Y_{i,t-1} + u_i + \varepsilon_{it}, \qquad \text{instrument } Y_{i,t-1} \text{ with } \Delta Y_{i,t-1}
\]

\begin{itemize}
\tightlist
\item
  \textbf{Why it helps.} A level cannot predict a random walk's next
  step --- but \textbf{a firm that just rose is high}. The first-stage
  coefficient of \(Y_{i,t-1}\) on \(\Delta Y_{i,t-1}\) is exactly
  \(\tfrac12\) for every \(\rho\). It never dies.
\item
  \textbf{What it costs.} The levels error contains \(u_i\), so you need
  \(\mathbb{E}[\Delta Y_{i,t-1}\,u_i] = 0\): \textbf{a firm's changes
  must say nothing about its permanent level.} True in the two-firm
  picture, \emph{because both firms start at their own long-run levels}:
  \textbf{mean stationarity}.
\item
  \textbf{When to doubt it.} Firms observed from IPO start away from
  their targets and converge; a panel that starts after a crisis has
  everyone off target together. Then rising \emph{does} reveal a
  high-\(u_i\) firm, and the extra instruments are invalid.
\end{itemize}
\end{frame}

\begin{frame}{Every estimator, on one line}
\protect\phantomsection\label{every-estimator-on-one-line}
\begin{center}
\includegraphics[width=0.92\linewidth,height=\textheight,keepaspectratio]{06b_dynamic_panel_files/figure-beamer/one-line-1.pdf}
\end{center}

\footnotesize Shaded: the bracket. Dashed: the truth. Dots: medians;
bars: 10th--90th percentiles over 200 draws. \normalsize

\begin{itemize}
\tightlist
\item
  \textbf{\(\rho = 0.6\):} every IV estimator lands inside the bracket,
  near the truth.
\item
  \textbf{\(\rho = 0.95\):} the difference estimators collapse (weak
  instruments); system GMM holds (\(0.98\)).
\item
  \textbf{Firms start at the average level, not their own:}
  Arellano--Bond recovers (\(0.95\)); system GMM lands at \(1.22\),
  \textbf{above OLS} --- the bracket catches it. \textbf{Each estimator
  fails exactly when its own assumption does.}
\end{itemize}
\end{frame}

\begin{frame}[fragile]{Long panels: a different route}
\protect\phantomsection\label{long-panels-a-different-route}
When \(T\) is large the problem is milder, and a different tool is
available.

\begin{itemize}
\tightlist
\item
  \textbf{Bias-corrected FE.} The fixed-effects bias is \(O(1/T)\)
  \textbf{with a known formula}: estimate FE, subtract the bias at
  \(\hat\rho\), iterate (Kiviet 1995; Bruno 2005 for unbalanced panels;
  Stata's \texttt{xtlsdvc}). No instruments, no weak-instrument problem
  --- but the other regressors must be \textbf{strictly exogenous},
  which finance controls rarely are.
\item
  \textbf{Long differences} --- \(Y_{it} - Y_{i,t-k}\) --- keep the
  instruments strong even for persistent series (Hahn, Hausman and
  Kuersteiner 2007).
\item
  \textbf{Horse race.} On finance-like panels --- short, persistent,
  unbalanced --- Flannery and Hankins (2013) find bias-corrected FE and
  system GMM generally reliable, and difference GMM frequently not.
\item
  \textbf{Rule of thumb.} \(T \le 10\), \(\rho\) moderate:
  Arellano--Bond with the diagnostics. \(\rho\) near one: system GMM,
  with mean stationarity defended. \(T \ge 20\) and exogenous controls:
  bias-corrected FE.
\end{itemize}
\end{frame}

\section{5. In Practice}\label{in-practice}

\begin{frame}{A case study: how fast do firms adjust leverage?}
\protect\phantomsection\label{a-case-study-how-fast-do-firms-adjust-leverage}
The partial-adjustment model is our equation with \(Y\) = leverage and
\(1-\rho\) the \textbf{speed of adjustment}.

\begin{center}\small
\begin{tabular}{llcc}
\toprule
Paper & Estimator & $\hat\rho$ & speed $1-\hat\rho$ \\
\midrule
Fama and French (2002) & pooled (Fama--MacBeth) & $\approx 0.85$ & $\approx 7$--$18\%$ a year \\
Flannery and Rangan (2006) & fixed effects (IV for lagged leverage) & $\approx 0.66$ & $\approx 34\%$ \\
Lemmon, Roberts and Zender (2008) & system GMM & $\approx 0.75$ & $\approx 25\%$ \\
Huang and Ritter (2009) & long differencing & $\approx 0.80$ & $\approx 17$--$23\%$ \\
\bottomrule
\end{tabular}
\end{center}

\begin{itemize}
\tightlist
\item
  \textbf{The ordering is the bracket.} Fixed effects reports the
  \emph{fastest} adjustment, pooled estimation the \emph{slowest}; the
  IV-type estimates sit between.
\item
  Whether firms rebalance toward targets or just drift \textbf{turned on
  the bias this lecture is about}: a half-life of \(1.7\) years and one
  of \(5\) are different theories of the firm.
\item
  Caveat: the papers differ in samples, periods and leverage
  definitions. The ordering is suggestive, not a controlled horse race.
\end{itemize}
\end{frame}

\begin{frame}[fragile]{Try it: simulate, estimate, grade}
\protect\phantomsection\label{try-it-simulate-estimate-grade}
\begin{columns}[T]
\begin{column}{0.68\linewidth}
\begin{Shaded}
\begin{Highlighting}[]
\KeywordTok{clear} \OtherTok{all}
\KeywordTok{set} \DecValTok{seed}\NormalTok{ 20260202}
\KeywordTok{local}\NormalTok{ rho = 0.6                      }\CommentTok{// then 0.95}
\KeywordTok{set} \KeywordTok{obs}\NormalTok{ 500}
\KeywordTok{gen}\NormalTok{ firm = }\DataTypeTok{\_n}
\KeywordTok{gen}\NormalTok{ u = rnormal()                    }\CommentTok{// firm effect}
\NormalTok{* }\BaseNTok{start}\NormalTok{ each firm }\FunctionTok{at}\NormalTok{ its own }\KeywordTok{long}\NormalTok{{-}}\KeywordTok{run} \DecValTok{level}
\KeywordTok{gen} \FunctionTok{y}\NormalTok{ = u/(1{-}}\OtherTok{\textasciigrave{}rho\textquotesingle{}}\NormalTok{) + rnormal(0, 1/}\FunctionTok{sqrt}\NormalTok{(1{-}}\OtherTok{\textasciigrave{}rho\textquotesingle{}}\NormalTok{\^{}2))}
\NormalTok{expand 8}
\KeywordTok{bysort}\NormalTok{ firm: }\KeywordTok{gen} \FunctionTok{year}\NormalTok{ = }\DataTypeTok{\_n}
\NormalTok{xtset firm }\FunctionTok{year}
\KeywordTok{by}\NormalTok{ firm: }\KeywordTok{replace} \FunctionTok{y}\NormalTok{ = }\OtherTok{\textasciigrave{}rho\textquotesingle{}}\NormalTok{*}\FunctionTok{y}\NormalTok{[}\DataTypeTok{\_n}\NormalTok{{-}1] + u + rnormal() }\KeywordTok{if} \DataTypeTok{\_n}\NormalTok{ \textgreater{} 1}

\KeywordTok{reg}   \FunctionTok{y}\NormalTok{ L.}\FunctionTok{y}\NormalTok{, }\KeywordTok{vce}\NormalTok{(}\KeywordTok{cluster}\NormalTok{ firm)                    }\CommentTok{// up}
\KeywordTok{xtreg} \FunctionTok{y}\NormalTok{ L.}\FunctionTok{y}\NormalTok{, }\KeywordTok{fe} \KeywordTok{vce}\NormalTok{(}\KeywordTok{cluster}\NormalTok{ firm)                 }\CommentTok{// down}
\NormalTok{ivregress 2sls D.}\FunctionTok{y}\NormalTok{ (LD.}\FunctionTok{y}\NormalTok{ = L2.}\FunctionTok{y}\NormalTok{), }\KeywordTok{vce}\NormalTok{(}\KeywordTok{cluster}\NormalTok{ firm)}
\KeywordTok{xtabond}  \FunctionTok{y}\NormalTok{, lags(1) maxldep(3) twostep }\KeywordTok{vce}\NormalTok{(}\KeywordTok{robust}\NormalTok{)}
\NormalTok{xtdpdsys }\FunctionTok{y}\NormalTok{, lags(1) maxldep(3) twostep }\KeywordTok{vce}\NormalTok{(}\KeywordTok{robust}\NormalTok{)}
\end{Highlighting}
\end{Shaded}
\end{column}

\begin{column}{0.3\linewidth}
\vspace{0.3em}
\begin{center}\scriptsize
One draw, $\hat\rho$:\\[0.3em]
\begin{tabular}{lccc}
\toprule
 & $0.6$ & $0.95$ & $0.95$\textsuperscript{off} \\
\midrule
OLS    & 0.918 & 0.999 & 1.083 \\
FE     & 0.334  & 0.592  & 0.840  \\
AH     & 0.596  & 1.077  & 1.004  \\
AB     & 0.609  & 0.824  & 0.966  \\
System & 0.589 & 1.010 & \textbf{1.150} \\
\bottomrule
\end{tabular}
\end{center}
\end{column}
\end{columns}

\footnotesize

Full file: \texttt{06b\_dynamic\_panel\_sim.do}. Set
\texttt{on\_target\ =\ 0} to start firms at the average level
(\textsuperscript{off}). One draw is noisy --- read it against the
figure, not on its own. \normalsize
\end{frame}

\begin{frame}[fragile]{The code, on real data}
\protect\phantomsection\label{the-code-on-real-data}
\begin{Shaded}
\begin{Highlighting}[]
\NormalTok{* 1. The bracket first {-}{-} and }\KeywordTok{keep}\NormalTok{ both numbers}
\KeywordTok{reg}   \FunctionTok{y}\NormalTok{ L.}\FunctionTok{y}\NormalTok{ x i.}\FunctionTok{year}\NormalTok{, }\KeywordTok{cluster}\NormalTok{(firm)}
\KeywordTok{xtreg} \FunctionTok{y}\NormalTok{ L.}\FunctionTok{y}\NormalTok{ x i.}\FunctionTok{year}\NormalTok{, }\KeywordTok{fe} \KeywordTok{cluster}\NormalTok{(firm)}

\NormalTok{* 2. Difference GMM (Arellano{-}Bond): }\KeywordTok{two}\NormalTok{{-}step, Windmeijer SEs, lags 2{-}4 }\KeywordTok{of} \FunctionTok{y}
\KeywordTok{xtabond} \FunctionTok{y}\NormalTok{ x i.}\FunctionTok{year}\NormalTok{, lags(1) maxldep(3) twostep }\KeywordTok{vce}\NormalTok{(}\KeywordTok{robust}\NormalTok{)}
\KeywordTok{estat}\NormalTok{ abond                       }\CommentTok{// AR(1) expected; AR(2) must be absent}

\NormalTok{* 3. Same, with collapsed instruments and the Hansen }\KeywordTok{test}\NormalTok{ (Roodman\textquotesingle{}s xtabond2)}
\KeywordTok{ssc}\NormalTok{ install xtabond2}
\NormalTok{xtabond2 }\FunctionTok{y}\NormalTok{ L.}\FunctionTok{y}\NormalTok{ x i.}\FunctionTok{year}\NormalTok{, gmm(}\FunctionTok{y}\NormalTok{, lag(2 4) }\KeywordTok{collapse}\NormalTok{) iv(x i.}\FunctionTok{year}\NormalTok{) noleveleq twostep }\KeywordTok{robust}
\NormalTok{*    x predetermined rather than exogenous? }\KeywordTok{replace}\NormalTok{ iv(x ...) with gmm(x, lag(1 3) }\KeywordTok{collapse}\NormalTok{)}

\NormalTok{* 4. System GMM {-}{-} only after you have defended }\KeywordTok{mean}\NormalTok{ stationarity}
\NormalTok{xtdpdsys }\FunctionTok{y}\NormalTok{ x i.}\FunctionTok{year}\NormalTok{, lags(1) maxldep(3) twostep }\KeywordTok{vce}\NormalTok{(}\KeywordTok{robust}\NormalTok{)}

\NormalTok{* 5. Long T, exogenous x: bias{-}corrected fixed }\KeywordTok{effects}
\KeywordTok{ssc}\NormalTok{ install xtlsdvc}
\NormalTok{xtlsdvc }\FunctionTok{y}\NormalTok{ x, initial(ab) bias(3) vcov(50)}
\end{Highlighting}
\end{Shaded}

\begin{itemize}
\tightlist
\item
  In R: \texttt{plm::pgmm()} with \texttt{transformation\ =\ "d"}
  (difference) or \texttt{"ld"} (system); \texttt{pdynmc} for the
  Roodman-style options.
\end{itemize}
\end{frame}

\begin{frame}{A reporting checklist}
\protect\phantomsection\label{a-reporting-checklist}
\begin{center}\small
\begin{tabular}{R{5.2cm}R{7.4cm}}
\toprule
Report & Because \\
\midrule
$\hat\rho_{\text{OLS}}$ and $\hat\rho_{\text{FE}}$ & the bracket; your estimate must sit inside it \\
AR(1) and AR(2) tests & AR(2) tests the identifying assumption \\
Hansen $J$, the instrument count, lag depth, collapsed or not & $J$ is informative only when the count is small relative to $N$ \\
one-step and two-step (Windmeijer) side by side & disagreement signals weakness \\
difference vs.\ system GMM, with mean stationarity defended & they rest on different assumptions \\
what $D$ is assumed to be (exogenous / predetermined / endogenous) & it determines which lags of $D$ are valid instruments \\
\bottomrule
\end{tabular}
\end{center}

\begin{itemize}
\tightlist
\item
  Each row is a place where the estimator can silently fail, and a
  number the reader needs to decide whether it did.
\end{itemize}
\end{frame}

\begin{frame}{Conclusion}
\protect\phantomsection\label{conclusion}
\begin{enumerate}
\tightlist
\item
  \textbf{A lagged \(Y\) is endogenous by construction}, not by
  omission. You can read the problem off the model.
\item
  \textbf{One rule does every calculation}: later shocks are unrelated
  to earlier \(Y\); the same-date shock enters with coefficient one.
  Check dates.
\item
  \textbf{OLS is biased up and fixed effects down.} Run both first ---
  \textbf{the bracket is the goalposts.}
\item
  \textbf{The fix is a longer lag of the same variable.} \(Y_{i,t-2}\)
  is one date too early to be contaminated. Arellano--Bond uses all such
  lags; cap and collapse them.
\item
  \textbf{The fix fails when the series is persistent or firm effects
  are large} --- the normal case in finance. Then: system GMM if firms
  sit at their own long-run levels, or bias-corrected FE if \(T\)
  allows.
\end{enumerate}

\vspace{0.3em}

\begin{itemize}
\tightlist
\item
  The question to carry out: \textbf{which way is this estimator wrong?}
  If you can say that, you can say where the truth lies.
\item
  \textbf{Next:} difference-in-differences --- when the switching has a
  \emph{date}.
\end{itemize}
\end{frame}

\begin{frame}{References}
\protect\phantomsection\label{references}
\tiny

\begin{itemize}
\tightlist
\item
  Anderson, T. W. and C. Hsiao (1982). Formulation and Estimation of
  Dynamic Models Using Panel Data. \emph{Journal of Econometrics} 18(1),
  47--82.
\item
  Angrist, J. and J.-S. Pischke (2009). \emph{Mostly Harmless
  Econometrics}, §§5.3--5.4 (bracketing lagged-dependent-variable and
  fixed-effects models).
\item
  Arellano, M. and S. Bond (1991). Some Tests of Specification for Panel
  Data: Monte Carlo Evidence and an Application to Employment Equations.
  \emph{Review of Economic Studies} 58(2), 277--297.
\item
  Blundell, R. and S. Bond (1998). Initial Conditions and Moment
  Restrictions in Dynamic Panel Data Models. \emph{Journal of
  Econometrics} 87(1), 115--143.
\item
  Bond, S. (2002). Dynamic Panel Data Models: A Guide to Micro Data
  Methods and Practice. \emph{Portuguese Economic Journal} 1(2),
  141--162.
\item
  Bruno, G. (2005). Approximating the Bias of the LSDV Estimator for
  Dynamic Unbalanced Panel Data Models. \emph{Economics Letters} 87(3),
  361--366.
\item
  Fama, E. and K. French (2002). Testing Trade-Off and Pecking Order
  Predictions about Dividends and Debt. \emph{RFS} 15(1), 1--33.
\item
  Flannery, M. and K. Hankins (2013). Estimating Dynamic Panel Models in
  Corporate Finance. \emph{Journal of Corporate Finance} 19, 1--19.
\item
  Flannery, M. and K. Rangan (2006). Partial Adjustment toward Target
  Capital Structures. \emph{JFE} 79(3), 469--506.
\item
  Hahn, J., J. Hausman and G. Kuersteiner (2007). Long Difference
  Instrumental Variables Estimation for Dynamic Panel Models with Fixed
  Effects. \emph{Journal of Econometrics} 140(2), 574--617.
\item
  Hansen, B. (2022). \emph{Econometrics}, ch.~17, §§17.36--17.40 --- the
  source for the \(T=3\) derivation and the weak-instrument formula.
\item
  Holtz-Eakin, D., W. Newey and H. Rosen (1988). Estimating Vector
  Autoregressions with Panel Data. \emph{Econometrica} 56(6),
  1371--1395.
\item
  Huang, R. and J. Ritter (2009). Testing Theories of Capital Structure
  and Estimating the Speed of Adjustment. \emph{JFQA} 44(2), 237--271.
\item
  Kiviet, J. (1995). On Bias, Inconsistency, and Efficiency of Various
  Estimators in Dynamic Panel Data Models. \emph{Journal of
  Econometrics} 68(1), 53--78.
\item
  Lemmon, M., M. Roberts and J. Zender (2008). Back to the Beginning:
  Persistence and the Cross-Section of Corporate Capital Structure.
  \emph{JF} 63(4), 1575--1608.
\item
  Nickell, S. (1981). Biases in Dynamic Models with Fixed Effects.
  \emph{Econometrica} 49(6), 1417--1426.
\item
  Roodman, D. (2009a). How to Do xtabond2: An Introduction to Difference
  and System GMM in Stata. \emph{Stata Journal} 9(1), 86--136.
\item
  Roodman, D. (2009b). A Note on the Theme of Too Many Instruments.
  \emph{Oxford Bulletin of Economics and Statistics} 71(1), 135--158.
\item
  Windmeijer, F. (2005). A Finite Sample Correction for the Variance of
  Linear Efficient Two-Step GMM Estimators. \emph{Journal of
  Econometrics} 126(1), 25--51.
\end{itemize}
\end{frame}

\section{Backup}\label{backup}

\begin{frame}{Bracketing \(\delta\): lagged \(Y\) or fixed effects?}
\protect\phantomsection\label{bracketing-delta-lagged-y-or-fixed-effects}
When you are unsure \textbf{which model} is right --- lagged \(Y\), or
fixed effects --- and what you care about is \(\delta\), the two
estimates bracket it too.

\begin{center}\small
\begin{tabular}{lll}
\toprule
If the truth is & and you estimate & then $\hat\delta$ is \\
\midrule
the lagged model & fixed effects      & too \textbf{high} \\
the fixed effects model & the lagged model & too \textbf{low} \\
\bottomrule
\end{tabular}
\end{center}

\begin{itemize}
\tightlist
\item
  As written, for \(\delta > 0\); if \(\delta < 0\) the directions swap.
  (Angrist and Pischke 2009, §§5.3--5.4.)
\item
  Run both and report the bracket rather than picking the specification
  you prefer.
\item
  \textbf{If the bracket excludes every sensible value, stop.} The model
  is wrong in some other way.
\end{itemize}
\end{frame}




\end{document}
